\documentclass[10pt,a4paper]{amsart}
\usepackage[margin=19mm]{geometry}
\usepackage{amsmath,amssymb,mathtools}
\usepackage[scr=rsfs,cal=euler]{mathalfa}
\usepackage{xfrac}
\usepackage[unicode,hidelinks,bookmarks=false]{hyperref}
\newcommand{\ie}{i.e.\ }
\newcommand{\cf}{cf.\ }
\newcommand{\resp}{respectively\ }
\newcommand{\Subsection}{\subsection}
\providecommand{\nobreakdash}{\nobreak\hbox{-}}
\setlength{\emergencystretch}{2em}
\multlinegap0pt
\usepackage{mathrsfs}
\usepackage{enumitem}
\hypersetup{
	colorlinks=true,
	linkcolor=blue,
	citecolor=blue,
	urlcolor=blue
}
\newtheorem{theorem}{Theorem}[section]
\newtheorem{lemma}[theorem]{Lemma}
\newtheorem{proposition}[theorem]{Proposition}
\newtheorem{corollary}[theorem]{Corollary}
\newtheorem{definition}[theorem]{Definition}
\newtheorem{remark}[theorem]{Remark}
\newcommand{\PSH}{\mathrm{PSH}}
\newcommand{\SHm}{\mathrm{SH}_m}
\newcommand{\ddc}{dd^c}
\newcommand{\vol}{\omega^n}
\begin{document}
\title[Stability and strong convergence for complex Hessian equatio]{Stability and strong convergence for complex Hessian equations with $L^1$ data}
\author{Truong Dinh Dat}
\date{}
\maketitle
\noindent\emph{Source and licence.} Stability and strong convergence for complex Hessian equations with $L^1$ data. Version arXiv:2607.04704v1 (CC BY 4.0). This material is distributed under the Creative Commons Attribution 4.0 International licence, http://creativecommons.org/licenses/by/4.0/, as recorded in the arXiv OAI licence metadata for this exact version and shown on the abstract page https://arxiv.org/abs/2607.04704v1. Full rights notice: licenses/arxiv-2607.04704-CC-BY-4.0.txt.\par\smallskip
\noindent\textit{Editorial scope.} Counted unit: the original \texttt{theorem} environment \texttt{thm:measure\_stability} at source lines 839--857 together with its complete proof at lines 859--961; the delivered document keeps source lines 291--962 so that every referenced label and every citation resolves inside it. Native editable source is the paper's own concatenated TeX; the AMS wrapper is editorial formatting only. No publisher class, upstream script or unrelated artwork is redistributed.
	\begin{lemma}\label{lem:AC}
		Let $u\in\mathcal F_m(\Omega)$. Then for any Borel set $E\subset\Omega$ and any $t>0$,
		\[
		H_m(u)(E)
		\le
		\int_{\{u<-t\}} H_m(u)
		+
		t^m\,\mathrm{Cap}_m(E,\Omega).
		\]
		In particular, $H_m(u)$ is absolutely continuous with respect to the
		$m$-Hessian capacity.
	\end{lemma}
	
	\begin{proof}
For completeness, we sketch the proof.
		Fix $t>0$ and let $E\subset\Omega$ be a Borel set.
		We decompose
		\[
		H_m(u)(E)
		=
		\int_{E\cap\{u<-t\}} H_m(u)
		+
		\int_{E\cap\{u\ge -t\}} H_m(u).
		\]
		
		The first term is trivially bounded by
		\[
		\int_{E\cap\{u<-t\}} H_m(u)
		\le
		\int_{\{u<-t\}} H_m(u).
		\]
		
		To estimate the second term, we consider the truncated function
		\[
		u_t := \max(u,-t),
		\]
		which belongs to $\mathcal F_m(\Omega)$ and satisfies
		\[
		- t \le u_t \le 0.
		\]
		Hence the function $u_t/t$ is $m$-subharmonic on $\Omega$ and satisfies
		\[
		-1 \le \frac{u_t}{t} \le 0.
		\]
		
		By the locality property of the Hessian operator,
		we have $H_m(u_t) = H_m(u)$ on the set $\{u>-t\}$ and since $(\ddc(u_t/t))^m = \frac{1}{t^m}(\ddc u_t)^m$. Therefore,
		\[
		\int_{E\cap\{u\ge -t\}} H_m(u)
		=
		\int_{E\cap\{u\ge -t\}} H_m(u_t)
		=
		t^m \int_{E\cap\{u\ge -t\}} H_m\!\left(\frac{u_t}{t}\right).
		\]
		
		By the definition of the $m$-Hessian capacity and the fact that
		$\frac{u_t}{t}$ is an admissible test function, we obtain
		\[
		\int_{E\cap\{u\ge -t\}} H_m(u)
		\le
		t^m \, \mathrm{Cap}_m(E,\Omega).
		\]
		
		Combining the above estimates yields
		\[
		H_m(u)(E)
		\le
		\int_{\{u<-t\}} H_m(u)
		+
		t^m \, \mathrm{Cap}_m(E,\Omega),
		\]
		which proves the desired inequality.
		
		In particular, letting $t\to\infty$ shows that $H_m(u)$ is absolutely
		continuous with respect to the $m$-Hessian capacity.
	\end{proof}
	\begin{remark}
		Lemma~\ref{lem:AC} provides a quantitative form of the absolute continuity
		of the Hessian measure with respect to the $m$-Hessian capacity,
		which plays a crucial role in the proof of strong convergence results.
	\end{remark}
	
	\section{A priori estimates and approximation}
	
	In this section we establish uniform estimates and approximation results for solutions of the complex Hessian equation.
	These results form the technical core of the paper and allow us to treat right-hand sides belonging merely to $L^1(\Omega)$.
	
	%------------------------------------------------
	
	\subsection{Approximation of the right-hand side}
	
	Let $f \in L^1(\Omega)$, $f \ge 0$.
	We consider the standard truncation
	\[
	f_j := \min(f,j), \quad j \in \mathbb{N}.
	\]
	Then $\{f_j\}$ is an increasing sequence of bounded functions satisfying
	\[
	f_j \in L^\infty(\Omega), 
	\quad
	0 \le f_j \le f_{j+1},
	\quad
	f_j \longrightarrow f \ \text{in } L^1(\Omega).
	\]
	
	For each $j$, it is known that there exists a unique solution
	$u_j \in \mathcal{E}_m^0(\Omega)$ of the Dirichlet problem
	\begin{equation}\label{eq:approx}
		H_m(u_j) = f_j \, \beta^n.
	\end{equation}
	This follows from the results of Dinew--Ko\l odziej.
	
	The main difficulty is to obtain estimates on $\{u_j\}$ that are uniform in $j$ and depend only on $\|f\|_{L^1(\Omega)}$.
	
	%------------------------------------------------
	
	\subsection{Uniform mass control}
	
	We begin with a simple but fundamental observation.
	
	\begin{lemma}\label{lem:mass}
		For all $j \in \mathbb{N}$, one has
		\[
		\int_\Omega H_m(u_j)
		\le \int_\Omega f \, \beta^n.
		\]
	\end{lemma}
	
	\begin{proof}
		Since $H_m(u_j)=f_j\beta^n$, we have
		\[
		\int_\Omega H_m(u_j)
		=
		\int_\Omega f_j \, \beta^n
		\le
		\int_\Omega f \, \beta^n,
		\]
		because $0 \le f_j \le f$.
	\end{proof}
	

	
	%------------------------------------------------
	
	\subsection{Capacity estimates}
	
	Uniform control of sublevel sets will be crucial for compactness arguments.
	
	\begin{lemma}\label{lem:capacity}
		 For every $t>0$ and all $j \in \mathbb{N}$,
		\[
		\mathrm{Cap}_m\big(\{u_j < -t\}, \Omega\big)
		\le \frac{1}{t^m} \int_\Omega f \, \beta^n.
		\]
	\end{lemma}
	
	\begin{proof}
		By the capacity estimate recalled in Section~2 and Lemma~\ref{lem:mass}, we obtain
		\[
		\mathrm{Cap}_m(\{u_j<-t\},\Omega)
		\le \frac{1}{t^m} \int_\Omega H_m(u_j)
		\le \frac{1}{t^m} \int_\Omega f \, \beta^n.
		\]
		This proves the claim.
	\end{proof}
	
	In particular, the sequence $\{u_j\}$ is uniformly tight with respect to the $m$-Hessian capacity.
	
	%------------------------------------------------
	
	\subsection{Compactness and convergence in capacity}
	
	We now establish compactness of the approximating solutions.
	
	\begin{lemma}\label{lem:compact}
		Let $(u_j)\subset SH_m(\Omega)$ be uniformly bounded above and satisfy
		\[
		\sup_j \mathrm{Cap}_m(\{u_j<-t\},\Omega)\to 0 \quad \text{as } t\to\infty.
		\]
		Then there exists a subsequence converging in $m$-Hessian capacity
		to some $u\in SH_m(\Omega)$.
	\end{lemma}
	
	\begin{proof}
		Since the sequence $(u_j)$ is uniformly bounded above and satisfies
		a uniform $m$-Hessian capacity estimate on sublevel sets
		(Lemma~\ref{lem:capacity}), this follows from the compactness theorem for bounded families
		of $m$-subharmonic functions with uniform capacity control
		(see e.g. \cite{DinewKolodziej2014}).
	\end{proof}
	
	
	Moreover, since $f_j$ is increasing, the comparison principle implies that $\{u_j\}$ is a decreasing sequence.
	Hence the whole sequence converges pointwise almost everywhere to $u$.
	
	%------------------------------------------------
	
	\subsection{Passage to the limit in the Hessian equation}
	
	We conclude this section by showing that the limit function solves the desired equation.
	
	\begin{lemma}\label{lem:limit}
		Let $u_j$ be the solutions of \eqref{eq:approx} and let $u$ be their limit.
		Then
		\[
		H_m(u) = f \, \beta^n
		\quad \text{in the sense of measures}.
		\]
	\end{lemma}
	
	\begin{proof}
		Since $(u_j)$ is a decreasing sequence in $\mathcal F_m(\Omega)$ with
		uniformly bounded Hessian mass, the monotone convergence theorem for
		the complex Hessian operator applies.
		\[
		H_m(u_j) \rightharpoonup H_m(u).
		\]
		On the other hand, $f_j \to f$ in $L^1(\Omega)$ implies
		\[
		f_j \, \beta^n \rightharpoonup f \, \beta^n
		\]
		weakly as measures.
		The claim follows.
	\end{proof}
	\begin{lemma}\label{lem:L1limit}
		Assume $f_j\to f$ in $L^1(\Omega)$.
		Then $f_j\beta^n$ converges weakly to $f\beta^n$.
	\end{lemma}
	\begin{proof}
		Since $f_j \to f$ in $L^1(\Omega)$, for any $\varphi \in C^0(\overline{\Omega})$ we have
		\[
		\int_\Omega \varphi f_j \,\beta^n \to \int_\Omega \varphi f \,\beta^n.
		\]
		Hence $f_j\beta^n$ converges weakly to $f\beta^n$.
	\end{proof}
	
	
	
	\section{Existence and stability of solutions}
	
	
	
	%------------------------------------------------
	
	\subsection{Existence theorem}
	
	\begin{theorem}\label{thm:existence}
		(see \cite{Lu})
		Let $\Omega \subset \mathbb{C}^n$ be a bounded $m$-hyperconvex domain and let
		$f \in L^1(\Omega)$, $f \ge 0$.
		Then there exists a function
		\[
		u \in \mathcal{F}_m(\Omega)
		\]
		such that
		\begin{equation}\label{eq:existence}
			(\ddc u)^m \wedge \beta^{n-m} = f \, \beta^n
		\end{equation}
		in the sense of measures.
		Moreover, the solution $u$ is unique in the class $\mathcal{F}_m(\Omega)$.
	\end{theorem}
	
	%------------------------------------------------
	
	\subsection{Proof of the existence}
	
	\begin{proof}
		Let $\{f_j\}$ be the truncation sequence defined by
		$f_j = \min(f,j)$.
		For each $j$, let $u_j \in \mathcal{E}_m^0(\Omega)$ be the unique solution of
		\[
		H_m(u_j) = f_j \, \beta^n.
		\]
		
		By Lemma~\ref{lem:mass}, the sequence $\{u_j\}$ has uniformly bounded Hessian mass.
		Lemma~\ref{lem:capacity} implies uniform control of sublevel sets in $m$-Hessian capacity.
		Hence, by Lemma~\ref{lem:compact}, there exists a subsequence converging in $m$-capacity to a function
		$u \in \mathcal{F}_m(\Omega)$.
		
		Since the sequence $\{f_j\}$ is increasing, the comparison principle yields that $\{u_j\}$ is decreasing.
		Therefore the whole sequence converges pointwise almost everywhere to $u$.
		By Lemma~\ref{lem:limit}, we obtain
		\[
		H_m(u) = f \, \beta^n,
		\]
		which proves the existence of a solution.
		
		Uniqueness follows from the comparison principle in the energy class $\mathcal{F}_m(\Omega)$.
	\end{proof}
	
	%------------------------------------------------
	
\begin{remark}
	The existence result follows from standard approximation arguments
	and is consistent with earlier works, in particular \cite{Lu}.
\end{remark}
\subsection{Stability in Hessian capacity}



\begin{theorem}\label{thm:stability}(see Theorem 7.2 \cite{V T Nguyen})
	Let $\{f_j\} \subset L^1(\Omega)$ be a sequence of nonnegative functions such that
	$f_j \to f$ in $L^1(\Omega)$.
	Let $u_j, u \in \mathcal{F}_m(\Omega)$ be the unique solutions of
	\[
	H_m(u_j) = f_j \, \beta^n,
	\qquad
	H_m(u) = f \, \beta^n.
	\]
	Then $u_j \to u$ in $m$-Hessian capacity.
\end{theorem}


\begin{proof}
	Let $\varepsilon>0$ be fixed. We will show that
	\[
	\mathrm{Cap}_m\big(\{|u_j-u|>\varepsilon\},\Omega\big)\to 0
	\quad \text{as } j\to\infty.
	\]
	
	We begin with the estimate of the set $\{u_j<u-\varepsilon\}$.
	
	By the comparison principle, we have
	\[
	\int_{\{u_j<u-\varepsilon\}} H_m(u_j)
	\le
	\int_{\{u_j<u-\varepsilon\}} H_m(u).
	\]
	Since $H_m(u_j)=f_j\,\beta^n$ and $H_m(u)=f\,\beta^n$, it follows that
	\[
	\int_{\{u_j<u-\varepsilon\}} f_j\,\beta^n
	\le
	\int_{\{u_j<u-\varepsilon\}} f\,\beta^n.
	\]
	Hence,
	\[
	\int_{\{u_j<u-\varepsilon\}} (f_j-f)\,\beta^n \le 0,
	\]
	which implies
	\[
	\int_{\{u_j<u-\varepsilon\}} f_j\,\beta^n
	\le
	\int_\Omega (f-f_j)_+\,\beta^n.
	\]
	
	Now, by the definition of the $m$-Hessian capacity, for any
	$v\in SH_m(\Omega)$ with $-1\le v\le 0$, we have
	\[
	\int_{\{u_j<u-\varepsilon\}} H_m(v)
	\le
	\frac{1}{\varepsilon^m}
	\int_{\{u_j<u-\varepsilon\}} H_m(u_j).
	\]
	Taking the supremum over all such $v$, we obtain
	\[
	\mathrm{Cap}_m\big(\{u_j<u-\varepsilon\},\Omega\big)
	\le
	\frac{1}{\varepsilon^m}
	\int_{\{u_j<u-\varepsilon\}} H_m(u_j).
	\]
	Combining the above estimates yields
	\[
	\mathrm{Cap}_m\big(\{u_j<u-\varepsilon\},\Omega\big)
	\le
	\frac{1}{\varepsilon^m}
	\int_\Omega (f-f_j)_+\,\beta^n.
	\]
	
	Since $f_j\to f$ in $L^1(\Omega)$, we have
	\[
	\int_\Omega (f-f_j)_+\,\beta^n \longrightarrow 0,
	\]
	hence
	\[
	\mathrm{Cap}_m\big(\{u_j<u-\varepsilon\},\Omega\big)\to 0.
	\]
	
	The estimate for the set $\{u<u_j-\varepsilon\}$ is obtained in the same way, by exchanging the roles of $u_j$ and $u$. Therefore,
	\[
	\mathrm{Cap}_m\big(\{u<u_j-\varepsilon\},\Omega\big)\to 0.
	\]
	
	Finally, since
	\[
	\{|u_j-u|>\varepsilon\}
	\subset
	\{u_j<u-\varepsilon\}\cup \{u<u_j-\varepsilon\},
	\]
	we conclude that
	\[
	\mathrm{Cap}_m\big(\{|u_j-u|>\varepsilon\},\Omega\big)\to 0.
	\]
	
	This proves that $u_j\to u$ in $m$-Hessian capacity.
\end{proof}
\begin{remark}
	The above proof avoids the use of non-admissible test functions 
	and relies only on the comparison principle and the $L^1$ convergence 
	of the densities.
\end{remark}
	\section{Strong convergence}
	
	In this section we establish a strong convergence result for solutions of the complex Hessian equation with $L^1$ right-hand side.
	This result constitutes the main contribution of the paper.	
	%------------------------------------------------
	

	
	
	
	%------------------------------------------------
	
	\subsection{Strong convergence in energy}
	
This result shows that the convergence holds in the natural
energy topology associated to the Hessian operator.
In particular, it implies that the nonlinear measures $H_m(u_j)$
interact well with the convergence of potentials,
which is not captured by capacity convergence alone.
	
	\begin{theorem}\label{thm:strong}
		Under the assumptions of Theorem~\ref{thm:stability}, the sequence $\{u_j\}$ converges strongly to $u$ in $\mathcal{F}_m(\Omega)$, namely
		\[
		\int_\Omega |u_j - u| \, H_m(u_j) \longrightarrow 0.
		\]
	\end{theorem}
	
\begin{proof}
	By Theorem~\ref{thm:stability}, we know that $u_j \to u$ in $m$-Hessian capacity. 
	Moreover, by Lemma~\ref{lem:mass}, the total masses $\int_\Omega H_m(u_j)$ are uniformly bounded.
	
	Let $\varepsilon>0$. We will show that
	\[
	\int_\Omega |u_j-u|\,H_m(u_j) < \varepsilon
	\]
	for all sufficiently large $j$.
	
	Fix $\delta>0$ (to be chosen later). We split
	\[
	\int_\Omega |u_j-u|\,H_m(u_j)
	\le 
	\delta \int_\Omega H_m(u_j)
	+
	\int_{\{|u_j-u|>\delta\}} H_m(u_j).
	\]
	
	Since $\sup_j \int_\Omega H_m(u_j) < +\infty$, we can choose $\delta>0$ such that
	\[
	\delta \sup_j \int_\Omega H_m(u_j) < \frac{\varepsilon}{3}.
	\]
	
	It remains to estimate the second term. 
	Applying Lemma~\ref{lem:AC} with $u = u_j$ and the Borel set 
	\[
	E_j := \{|u_j-u|>\delta\},
	\]
	we obtain for any $t>0$:
	\[
	H_m(u_j)(E_j)
	\le 
	\int_{\{u_j<-t\}} H_m(u_j)
	+
	t^m \,\mathrm{Cap}_m(E_j,\Omega).
	\]
	
	We now control the two terms on the right-hand side.
	
	\medskip
	\noindent
	\textbf{Step 1. Control of the sublevel term.}
	
	By Lemma~\ref{lem:capacity}, we have
	\[
	\int_{\{u_j<-t\}} H_m(u_j)
	\le 
	\frac{1}{t^m} \int_\Omega f\,\beta^n.
	\]
	Hence we can choose $t>0$ sufficiently large so that
	\[
	\sup_j \int_{\{u_j<-t\}} H_m(u_j) < \frac{\varepsilon}{3}.
	\]
	
	\medskip
	\noindent
	\textbf{Step 2. Control of the capacity term.}
	
	Since $u_j \to u$ in $m$-Hessian capacity, we have
	\[
	\mathrm{Cap}_m(E_j,\Omega)
	=
	\mathrm{Cap}_m(\{|u_j-u|>\delta\},\Omega)
	\longrightarrow 0.
	\]
	Therefore, for the above fixed $t$, there exists $j_0$ such that for all $j \ge j_0$,
	\[
	t^m \,\mathrm{Cap}_m(E_j,\Omega)
	<
	\frac{\varepsilon}{3}.
	\]
	
	\medskip
	Combining the above estimates, we obtain for all $j \ge j_0$:
	\[
	\int_\Omega |u_j-u|\,H_m(u_j)
	<
	\frac{\varepsilon}{3}
	+
	\frac{\varepsilon}{3}
	+
	\frac{\varepsilon}{3}
	=
	\varepsilon.
	\]
	
	This proves the desired convergence.
\end{proof}
\begin{remark}
	The strong convergence result in Theorem~\ref{thm:strong} should be compared 
	with the analogous theory for the complex Monge--Amp\`ere equation (the case $m=n$). 
	In that setting, strong stability results under $L^1$ convergence of the densities 
	have been established by several authors (see e.g. \cite{DinewPlis, Liu-Zhang}), 
	relying on the full strength of the comparison principle and the rich structure 
	of the Monge--Amp\`ere operator.
	
	In contrast, for complex Hessian equations with $1 \le m < n$, the situation is 
	more delicate. The Hessian operator lacks some of the key structural properties 
	available in the Monge--Amp\`ere case, such as full monotonicity and a stronger 
	comparison framework. As a consequence, several arguments used in the 
	Monge--Amp\`ere setting do not directly carry over to the Hessian context.
	
	The proof of Theorem~\ref{thm:strong} shows that, despite these difficulties, 
	one can still obtain strong convergence in the natural energy topology 
	by combining capacity convergence, uniform control of sublevel sets, 
	and a quantitative form of absolute continuity of Hessian measures 
	with respect to the $m$-Hessian capacity (Lemma~\ref{lem:AC}).
	
	This provides a nontrivial extension of known stability results to the 
	Hessian setting with merely $L^1$ data.
\end{remark}
\section{Application to measure data}

In this section we indicate how the strong convergence result obtained above
extends to the case of general measure data with uniformly bounded mass.
The argument follows the same strategy as in the $L^1$ case, combining
capacity estimates, compactness, and the absolute continuity of Hessian measures.

\medskip


\begin{theorem}\label{thm:measure_stability}
	Let $\{\mu_j\}$ be a sequence of nonnegative Radon measures on $\Omega$
	such that
	\[
	\sup_j \mu_j(\Omega) < +\infty,
	\]
	and assume that $\mu_j \rightharpoonup \mu$ weakly as measures.
	
	Let $u_j, u \in \mathcal{F}_m(\Omega)$ be the unique solutions of
	\[
	H_m(u_j) = \mu_j,
	\qquad
	H_m(u) = \mu.
	\]
	Then $u_j \to u$ in $m$-Hessian capacity. Moreover,
	\[
	\int_\Omega |u_j - u| \, H_m(u_j) \longrightarrow 0.
	\]
\end{theorem}

\begin{proof}
	We divide the proof into two steps.
	
	\medskip
	\noindent
	\textbf{Step 1. Convergence in capacity.}
	
	Since $H_m(u_j)=\mu_j$ and $\sup_j \mu_j(\Omega)<+\infty$, we have
	uniform control of the Hessian mass:
	\[
	\int_\Omega H_m(u_j) = \mu_j(\Omega) \le C.
	\]
	By the capacity estimate recalled in Section~2, for every $t>0$,
	\[
	\mathrm{Cap}_m(\{u_j<-t\},\Omega)
	\le \frac{C}{t^m}.
	\]
	Thus the sequence $\{u_j\}$ is uniformly tight with respect to the
	$m$-Hessian capacity.
	
	By the compactness theorem for $m$-subharmonic functions with uniform
	capacity control (see e.g. \cite{DinewKolodziej2014}), there exists a subsequence
	converging in $m$-Hessian capacity to some $v \in \mathcal{F}_m(\Omega)$.
	
	On the other hand, since $H_m(u_j)=\mu_j$ and $\mu_j \rightharpoonup \mu$,
	it follows that any such limit $v$ satisfies
	\[
	H_m(v) = \mu
	\]
	in the sense of measures (see e.g. \cite{DinewKolodziej2014, Lu}).
	By uniqueness of solutions in $\mathcal{F}_m(\Omega)$, we conclude that $v=u$.
	Therefore, the whole sequence $u_j \to u$ in $m$-Hessian capacity.
	
	\medskip
	\noindent
	\textbf{Step 2. Strong convergence in energy.}
	
	We now prove that
	\[
	\int_\Omega |u_j-u|\,H_m(u_j) \longrightarrow 0.
	\]
	
	Let $\varepsilon>0$. As in the proof of Theorem~\ref{thm:strong}, we write
	\[
	\int_\Omega |u_j-u|\,H_m(u_j)
	\le
	\delta \int_\Omega H_m(u_j)
	+
	\int_{\{|u_j-u|>\delta\}} H_m(u_j).
	\]
	
	Since $\int_\Omega H_m(u_j)\le C$, we can choose $\delta>0$ such that
	\[
	\delta C < \frac{\varepsilon}{3}.
	\]
	
	For the second term, applying Lemma~\ref{lem:AC} with $u=u_j$ and
	$E_j := \{|u_j-u|>\delta\}$, we obtain for any $t>0$:
	\[
	H_m(u_j)(E_j)
	\le
	\int_{\{u_j<-t\}} H_m(u_j)
	+
	t^m\,\mathrm{Cap}_m(E_j,\Omega).
	\]
	
	By the capacity estimate, we have
	\[
	\int_{\{u_j<-t\}} H_m(u_j)
	=
	\mu_j(\{u_j<-t\})
	\le \frac{C}{t^m}.
	\]
	Hence we can choose $t>0$ sufficiently large so that
	\[
	\sup_j \int_{\{u_j<-t\}} H_m(u_j)
	<
	\frac{\varepsilon}{3}.
	\]
	
	On the other hand, since $u_j \to u$ in $m$-Hessian capacity,
	\[
	\mathrm{Cap}_m(E_j,\Omega)
	=
	\mathrm{Cap}_m(\{|u_j-u|>\delta\},\Omega)
	\longrightarrow 0.
	\]
	Thus, for $j$ sufficiently large,
	\[
	t^m\,\mathrm{Cap}_m(E_j,\Omega)
	<
	\frac{\varepsilon}{3}.
	\]
	
	Combining the above estimates, we obtain for all large $j$:
	\[
	\int_\Omega |u_j-u|\,H_m(u_j)
	<
	\varepsilon.
	\]
	
	This proves the desired convergence.
\end{proof}


\begin{thebibliography}{99}
\bibitem[DinewKolodziej2014]{DinewKolodziej2014} S. Dinew and S. Ko{\l}odziej, A priori estimates for complex Hessian equations, \emph{Anal. PDE} \textbf{7} (2014), 227--244.
\bibitem[DinewPlis]{DinewPlis} S. Dinew and S. P{\l}is, Uniqueness in $L^1$ for complex Monge--Amp\`ere equations, \emph{Math. Z.} \textbf{277} (2014), 937--952.
\bibitem[Liu-Zhang]{Liu-Zhang} S. Liu and L. Zhang, $L^1$ stability for complex Monge--Amp\`ere equations, preprint, arXiv:2503.18392.
\bibitem[Lu]{Lu} Chinh H. Lu, A variational approach to complex Hessian equations in $\mathbb{C}^n$, \emph{J. Math. Anal. Appl.} \textbf{431} (2015), 228-259.
\bibitem[V T Nguyen]{V T Nguyen} Van Thien Nguyen, Maximal m-subharmonic functions and the Cegrell class $\mathcal{N}_m$, \emph{Indagationes Mathematicae} \textbf{30} (2019), 717-739.
\end{thebibliography}
\end{document}
